\chapter{Conclusion}
\label{ch:conc}

\section{Summary}
\label{sec:conc-summary}

This project will build and evaluate a CNN application that flags possible abnormalities in ECG images, reports calibrated confidence, and explains its predictions to different users. It will train on images rendered from PTB-XL and test on real images, addressing four gaps in the literature: limited testing on real images, missing calibration, unevaluated explanations, and little evidence on presenting results to different users.

The work is organised into seven phases (Section~\ref{sec:dev-plan}), with requirements tested through the Requirements Traceability Matrix and the project evaluated against its objectives, including a user study under the approved ethics application. Two prototypes, rendering and preprocessing, are complete, and GUI wireframes are planned. The main risks are the rendered-to-real gap, unreliable detection on poor images, and rare classes, and each has a fallback in the plan (Appendix~\ref{app:risks}).

\section{Expected Contributions}
\label{sec:contributions}

The project is expected to contribute three things beyond a working application:
\begin{enumerate}\setlength{\itemsep}{0pt}\setlength{\parskip}{0pt}
  \item \textbf{A per-class measure of the rendered-to-real gap.} Training on rendered images and testing on real ones of the same recordings, with leakage removed, will show which classes transfer to real images and which do not, and whether timing features reduce the gap (SQ1).
  \item \textbf{Calibration under data shift for a multi-label image classifier.} Measuring the calibration error per class before and after temperature scaling, on both rendered and real images, will show whether confidence scores fitted on rendered images can still be trusted on real ones (SQ2).
  \item \textbf{Evidence on presenting results to non-expert users.} The user study will show whether students can interpret calibrated confidence and Grad-CAM heatmaps, which few studies have evaluated with users (SQ3).
\end{enumerate}

\section{Next Steps}
\label{sec:next-steps}
 Semester 2 will follow the gantt schedule (Table~\ref{fig:gantt-s2}) in Appendix~\ref{app:gantt}. The first few weeks will render the full training set, match the real images to their labels and remove their patients from the training folds, and build the detection step, as every later phase depends on these. CNN training will then follow, with the GUI developed in parallel to it, aiming for a working prototype by week 9. The user study and evaluation will take place in weeks 9--11, leaving the final weeks for writing the dissertation. If the schedule slips, the fallbacks in Section~\ref{sec:dev-plan} will be applied in order, starting with dropping the optional features of Phase 6.